\paragraph{Incremental computation.}
Incremental computing is a paradigm for computing based on the idea of saving time by leveraging previous computational results.
It has a long history, going back to the 1964 and 1967 papers of Lombardi \cite{LISP:LR64,AIC:Lombardi67}.
There have been contributions from many fields, including theory
(dynamic graph algorithms \cite{ausiello1990dynamic,ramalingam1996incremental} and dynamic data structures \cite{bentley1980decomposable,leeuwen1980dynamization}), databases (view maintenance \cite{Book:GM99,chirkova2012materialized}), programming languages (incremental lambda-calculus \cite{field1990incremental}, incremental parsing \cite{ghezzi1979incremental}, interactive program-development environments \cite{reps1983incremental}, incremental dataflow analysis \cite{ryder1988incremental}, differential propagation in iterative fixpoint solvers \cite{fecht1998propagating}, and finite-differencing transformations in optimizing compilers \cite{paige1982finite}), document-preparation systems \cite{chen2000incremental}, constraint solvers \cite{freeman1990incremental}, truth-maintenance systems \cite{doyle1979truth}, and no doubt many others.
For a recent survey of incremental computation, see \cite{DBLP:conf/pepm/Liu24}.

Some of the work described above is only for limited models of computation.
DBSP is a Turing-complete language.
Because of its generality, DBSP could be applied to many of the application areas addressed by previous work, although it would not necessarily provide the same runtime guarantees of specialized results.
However, many of these applications would need the ability to detect fixpoints of nested computations.
The work described in this paper contributes to realizing the potential of DBSP by addressing an important semantic lacuna in the DBSP formalism.

More specifically, prior work has studied incremental recursive and fixpoint computation from complementary perspectives. Alvarez-Picallo et al.~\cite{AlvarezPicallo2019Fixing} develop a general theory of derivatives of fixpoints for incremental evaluation of recursive computations and Datalog. Abo Khamis et al.~\cite{AboKhamis2024Convergence} characterize algebraic conditions under which Datalog over pre-semirings converges in finitely many iterations and semi-naive evaluation applies. These works provide incremental fixpoint constructions or convergence characterizations for particular Datalog semantics. Our contribution is complementary: we formalize FPD as a finite-prefix runtime problem for incrementally transformed circuits and establish its general limitations and its exact detectability for useful circuit classes.

%
%
%
%
%
%

%
%
%
%
%



\paragraph{Feldera implementation of DBSP}
As mentioned in \secref{sec:intfp} and \secref{sec:detector}, the Feldera implementation of DBSP~\cite{feldera_repo} uses internal-state stability to detect convergence. Its runtime gives each operator a scope-indexed fixed-point check, intended to report whether the operator's output will remain stable under stable inputs, and declares a circuit stable when all of its operators report stability. Our theory gives this established implementation strategy a declarative account: IntFP and IntConv specify internal-state stability denotationally, StFP connects the declarative criterion to circuit state, and the StFP detector gives an algorithmic formalization.

The main novelty lies in the formal guarantees, especially the completeness result. We prove that IntConv is sound and that the StFP detector is sound and complete. The global IntConv detector combines the StFP detector with fixed-input and zero-output checks across all bracketed sub-circuits, yielding sound and complete IntConv detection. We also prove that IntConv is complete for a large class of useful circuits, including regular circuits and their incrementally optimized forms (\secref{sec:useful}). Thus, for these circuits, the internal-state-stability criterion detects every convergence admitted by the external specification. These results are machine-checked in Lean 4. Moreover, our level-indexed predicates and detectors give a formal, compositional account of nested bracketed circuits, whereas the Feldera SQL-to-DBSP compiler~\cite{feldera_repo} rejects a nested recursive operator within recursive code. Therefore, our theory may help guide future extended implementations.

\paragraph{Differential Dataflow and its mathematical foundation.}

Differential Dataflow (DD)~\cite{McSherry2013DifferentialD} also incrementalizes recursive queries over versions $T \times \N$, the DD counterpart of the delta-of-deltas setting in our counterexample. At the collection level, \texttt{Egress} selects the first repeated collection, at $i^* = \min\{i : W_{(t,i)} = W_{(t,i-1)}\}$; its differential implementation instead operates on $\delta W_{(t,i)}$. In this setting, $\delta W_{(t,i)}=0$ does not imply a repetition, since $W_{(t,i)}-W_{(t,i-1)}=\sum_{s\leq t}\delta W_{(s,i)}$. Thus, FirstZero cannot soundly determine $i^*$, and the paper defines differential \texttt{Egress} using $i^*$ without giving a finite-prefix criterion for obtaining it from the differential trace. Similar to DBSP, this also turns out to be a gap between the theory and the implementation of DD: its prototype instead reports convergence when no outstanding differences remain in the dataflow, a global internal-quiescence check similar in spirit to IntConv.

A later developed mathematical foundation for DD by Abadi et al.~\cite{Abadi2015FoundationsOD}, distinguishes fixed-iteration egress from first-repetition egress, but formalizes only the fixed-iteration policy; it notes that first-repetition egress would require partial streams. Consequently, its correctness theorem does not provide a compositional denotational account of data-dependent first-repetition termination for unbounded incremental iteration. Similarly, \textit{Flo}~\cite{Laddad2024FloAS}, a semantic foundation for progressive stream processing, requires bounded inner streams, leaving termination of unbounded inner iteration outside its scope. Our theory instead gives a sound criterion for such unbounded computations and identifies useful circuit classes for which that criterion is complete.
